\section{Introduction}
A dense statevector simulation of $n$ qubits stores $2^n$ complex amplitudes~\cite{TODO-statevector-memory}. On a microcontroller the state competes with the program and the runtime for kilobytes of RAM, and the arithmetic runs on 8-bit or soft-float cores. Whether a given circuit can be simulated on such a device, and whether it is better to send it to a PC, is a practical question for embedded and educational quantum tooling~\cite{TODO-embedded-quantum-sim}.

We ask: \emph{what limits statevector circuit simulation on microcontrollers (Arduino \uno, \espA, \espB), how accurately can those limits be predicted, and where is the local-versus-offload trade-off?} This is a characterisation study: the analysis rules were refined while the data were being examined, and every such choice is disclosed in Sections~\ref{sec:method} and~\ref{sec:limits}; we do not present it as a pre-registered hypothesis test.

Contributions: (i) the memory and time limit of each board and circuit, with the evidence type for each limit; (ii) a prediction of those limits: a memory rule, and two time models (M1: one constant per board, M2: one cost per gate kind) with held-out scores, identifiability checks and a sensitivity analysis; (iii) the local-versus-offload crossover as brackets of measured $n$ under the fastest and slowest PC simulation speed we observed, so the host PC is part of the result.
